\chapter{计算性能}

前向与反向都是计算图上的算子序列。性能取决于图如何编译、设备如何重叠计算与通信、以及参数如何在多卡多机上同步。
命令式便于写控制流，符号式便于整图优化。下面把二者写成同一套映射。

\section{编译器和解释器}

命令式编程按语句改状态。Python 解释执行，每步都经过前端。

\subsection{符号式编程}

先定义计算流程，再编译成可执行程序，最后喂入输入。
记命令式求值为
\[
(a,b,c,d)
\;\longmapsto\;
e=a+b,\; f=c+d,\; g=e+f,
\]
符号式先构造图 $G$，再
\[
G\longmapsto P=\operatorname{compile}(G),\qquad
P(a,b,c,d)=g.
\]
编译器可见全图，可消去中间量，甚至把常量折叠成 $P\equiv 10$。

\subsection{混合式编程}

PyTorch 默认动态图。TorchScript 把已追踪的命令式模型编译成静态图，
开发仍用 Python，部署用 $P=\operatorname{script}(f)$。

\subsection{Sequential的混合式编程}

$L$ 层序列 $f=f_L\circ\cdots\circ f_1$。解释器对每层单独发指令；
多 GPU 时单线程前端成为瓶颈。
把 $\operatorname{Sequential}$ 脚本化后，整条复合映射一次下发。

\subsubsection{通过混合式编程加速}

同一输入 $\bm{x}$，比较 $f(\bm{x})$ 与 $\operatorname{script}(f)(\bm{x})$ 的墙钟时间。
编译后减少 Python 往返，输出值不变。

\subsubsection{序列化}

\[
(f,\bm{\theta})
\;\longmapsto\;
\operatorname{save}(P,\bm{\theta})
\;\longmapsto\;
\text{其他设备或语言上的 }P.
\]
序列化的是编译后的图与参数，不是 Python 源。

\subsection{小结}

命令式写 $f$；符号式把计算图 $G$ 编译成 $P$。
混合式在 $f\mapsto P$ 后加速，且 $P(\bm{x})=f(\bm{x})$。

\subsection{练习}

核验：对已有模型做 $f\mapsto\operatorname{script}(f)$ 后，前向时延是否下降、数值是否一致。

\section{异步计算}

GPU 核函数默认异步：前端插入任务后立即返回，后端在设备上执行。

\subsection{通过后端异步处理}

前端发出 $\bm{C}=\bm{A}\bm{B}$ 后，Python 已继续；实际 GEMM 在队列中。
同步 $S$ 强制等待队列清空。测时若不调用 $S$，测到的是入队时间而非计算时间。

\subsection{障碍器与阻塞器}

障碍 $S$ 使前端与指定设备上所有流对齐：
\[
\operatorname{sync}(\text{device}):
\quad
\text{等待该设备队列空}.
\]
把结果拷回 CPU 或打印标量也会隐式阻塞。

\subsection{改进计算}

设前端入队、后端计算、回传分别为 $t_1,t_2,t_3$。
同步执行 $N$ 次总时间为 $N(t_1+t_2+t_3)$；
异步且 $Nt_2$ 主导时约为 $t_1+N t_2+t_3$。
队列过长会堆积中间张量，故常按小批量同步一次。

\subsection{小结}

前端与后端解耦：命令快速入队，计算并行执行。
过度填充队列增加显存；每个小批量一次 $\operatorname{sync}$ 使两侧大致对齐。

\subsection{练习}

核验：同一矩阵乘在 CPU 上是否仍能观察到入队与完成之间的时间差。

\section{自动并行}

框架根据计算图上的依赖决定哪些算子可并行。
无依赖的子图可同时占用不同设备。

\subsection{基于GPU的并行计算}

在 GPU1、GPU2 上各做 $10$ 次矩阵乘。
两次 $\operatorname{sync}$ 串行测时约为两段之和；去掉中间同步后，
墙钟时间接近 $\max(T_1,T_2)$，即自动重叠。

\subsection{并行计算与通信}

计算 $\bm{y}=f(\bm{x})$ 于 GPU，再 $\bm{y}\mapsto\mathrm{CPU}$。
$\operatorname{copy}(\texttt{non\_blocking})$ 允许尚未完成的后续核与总线传输重叠。
反传中先就绪的梯度可提前走 PCI-Express。

\subsection{小结}

计算图给出依赖；无依赖的计算与通信可并行。
设备间搬移与算子执行重叠，才能吃满总线与计算单元。

\subsection{练习}

核验：无依赖的多个小算子是否被自动并行，以及 CPU–GPU 同时计算加拷贝的总时延。

\section{硬件}

算法决定统计效率；硬件决定一次 $\nabla f$ 要多久。
延迟差几个数量级就会使可训练与不可训练对换。

\subsection{计算机}

一台训练机构成映射
\[
(\mathrm{CPU},\;\mathrm{RAM},\;\mathrm{GPU},\;\mathrm{PCIe},\;\mathrm{NIC},\;\text{存储})
\;\longmapsto\;
\text{每秒可完成的迭代}.
\]
CPU 跑调度与数据管道；参数与激活放 RAM/显存；PCIe 接加速卡。

\subsection{内存}

带宽与延迟同时约束。DDR 通道峰值约 $40$–$100\,\mathrm{GB/s}$，
但先发地址再突发读；小块随机访问远低于峰值。
深度学习应尽量连续、大批量访问。

\subsection{存储器}

持久存储同样由带宽与 IOPS 刻画，器件间差几个数量级。

\subsubsection{硬盘驱动器}

约 $7200\,\mathrm{RPM}$，最坏旋转等待约 $8\,\mathrm{ms}$，约 $10^2$ IOPS。
顺序带宽远高于随机。

\subsubsection{固态驱动器}

$10^5$–$5\times 10^5$ 量级 IOPS，带宽 $1$–$3\,\mathrm{GB/s}$。
按块擦写，随机小写远慢于读。

\subsubsection{云存储}

IOPS 与容量可配置。训练若大量小记录，延迟高时应提高 IOPS 配额。

\subsection{CPU}

核心执行指令，缓存降低访存延迟，向量单元做 SIMD。

\subsubsection{微体系结构}

前端取指、预测、解码；执行端口可同时发射多条独立微操作。
整数端口与浮点端口不同，吞吐取决于指令能否并行。

\subsubsection{矢量化}

SIMD 在一个周期内对向量寄存器做同一运算。
FMA 计算 $a\leftarrow a+b\cdot c$。寄存器可达 $512$ 位。

\subsubsection{缓存}

$4$ 核、$2\,\mathrm{GHz}$、每周期 $256$ 位 AVX 约需
\[
4\times 256\,\mathrm{bit}\times 2\times 10^9/\mathrm{s}=256\,\mathrm{GB/s},
\]
而内存接口只有数十 $\mathrm{GB/s}$。差额必须由 L1/L2/L3 命中补上。

\subsection{GPU和其他加速卡}

训练要存中间激活，常用 FP16/FP32；推断可 INT8 且不必存反传状态。
GPU 以吞吐换取相对 CPU 高一个数量级的密集线性代数。

\subsection{网络和总线}

PCIe 单链路高带宽、微秒级延迟，通道数有限。
NVLink 卡间带宽高于 PCIe；以太网跨机约 $10\,\mathrm{GB/s}$ 量级，比 NVLink 慢约一个数量级。

\subsection{更多延迟}

缓存命中约纳秒，内存约几十纳秒，PCIe/网络微秒到毫秒，磁盘毫秒。
同一迭代里应避免把计算拆成大量高延迟小传输。

\subsection{小结}

量少次多的传输浪费延迟。向量化与对齐决定能否达到峰值。
训练精度过低会下溢；算法应让热参数留在缓存带宽内。

\subsection{练习}

核验：对齐与非对齐访问、固定步幅访问的带宽，以及 FP16 累积是否溢出。

\section{多GPU训练}

$k$ 张 GPU 增加显存与算力。需把一次小批量上的 $\nabla f$ 拆开再聚合。

\subsection{问题拆分}

层并行按层切网络，模型并行按层内切参数，数据并行按样本切批量。
前两者要编排激活与梯度的传递；数据并行只同步梯度，实现最直接。

\subsection{数据并行性}

$k$ 张卡各持一份完整 $\bm{\theta}$。一次迭代：
\begin{align*}
\mathcal{B}&=\mathcal{B}_1\cup\cdots\cup\mathcal{B}_k,
\\
\bm{g}^{(j)}&=\nabla_{\bm{\theta}}\frac{1}{|\mathcal{B}_j|}
\sum_{i\in\mathcal{B}_j}\ell_i,
\\
\bm{g}&=\operatorname{allreduce}(\bm{g}^{(1)},\ldots,\bm{g}^{(k)}),
\\
\bm{\theta}&\leftarrow\bm{\theta}-\eta\bm{g}
\quad\text{在每张卡上相同}.
\end{align*}

\subsection{简单网络}

用修改后的 LeNet 作为 $f_{\bm{\theta}}$，以便显式写出参数拷贝与 $\operatorname{allreduce}$。

\subsection{数据同步}

$\operatorname{get\_params}$ 把 $\bm{\theta}$ 复制到设备 $j$ 并附加梯度缓冲。
$\operatorname{allreduce}$ 实现
\[
\bm{g}\leftarrow\sum_{j=1}^{k}\bm{g}^{(j)},
\]
再广播回每张卡。

\subsection{数据分发}

\[
(\bm{X},\bm{y})
\;\longmapsto\;
\bigl\{(\bm{X}^{(j)},\bm{y}^{(j)})\bigr\}_{j=1}^{k},
\qquad
\sum_j|\mathcal{B}_j|=|\mathcal{B}|.
\]

\subsection{训练}

每个小批量：分发数据、各卡前向反向、$\operatorname{allreduce}$、本地更新。
有效批量变为 $k|\mathcal{B}_j|$ 时，$\eta$ 通常需略增大。
单卡精度评估可只在一张卡上进行。

\subsection{小结}

数据并行把 $\mathcal{B}$ 切到 $k$ 张卡，聚合梯度后再广播。
更大有效批量需要相应调整 $\eta$。

\subsection{练习}

核验：批量从 $b$ 扩到 $kb$ 时精度与 $\eta$ 应如何同扩，以及分块 $\operatorname{allreduce}$ 的通信量。

\section{多GPU的简洁实现}

高级 API 把上一节的分发、反传与聚合收成同一训练循环。至少两张 GPU。

\subsection{简单网络}

用修改后的 ResNet-18：更小卷积核与步幅，去掉开头最大汇聚，以适应小图。

\subsection{网络初始化}

必须在目标设备上初始化 $\bm{\theta}$，再访问该设备上的参数，否则引用落在错误设备。

\subsection{训练}

循环执行：全设备初始化；把 $\mathcal{B}$ 切到各设备；并行求损失与梯度；聚合后更新；并行算精度。
当计算时间明显长于同步时间，加速比接近线性。

\subsection{小结}

单卡可自动求 $f_{\bm{\theta}}(\bm{x})$。
多卡时每设备先初始化，优化器自动 $\operatorname{allreduce}$ 梯度。

\subsection{练习}

核验：更多卡、更大 $b$ 与 $\eta$ 的组合，以及 CPU 与 GPU 算力不均时如何划分前后向。

\section{参数服务器}

从单机多卡到多机时，互连带宽从 NVLink、PCIe 降到以太网，同步算法必须适配拓扑。

\subsection{数据并行训练}

各 GPU 算 $\bm{g}_{ijk}$，先机内聚合，再跨机聚合，最后广播新 $\bm{\theta}$。
聚合位置可以是某张 GPU、CPU，或按参数块分散到不同设备。

\subsection{环同步（Ring Synchronization）}

$k$ 张卡组成环，每步向邻居发送 $1/k$ 的梯度块并接收一块，共 $2(k-1)$ 步后每卡得到全和。
在 NVLink 带宽远高于 PCIe 的机器上，环同步走 NVLink，避免单点经 PCIe 汇聚。

\subsection{多机训练}

每台机器：读不同批量、卡间切分、本地前向反向、机内聚合，再把梯度发给参数服务器（或对等节点），取回更新后的 $\bm{\theta}$。
机器间时钟差要求显式同步，否则同步 SGD 会等待最慢的那台。

\subsection{键值存储}

参数服务器把切片 $i$ 上的梯度写成
\[
\bm{g}_{i}
=
\sum_{k\in\mathrm{workers}}
\sum_{j\in\mathrm{GPUs}}
\bm{g}_{ijk}.
\]
$\operatorname{push}$ 写入 $\bm{g}_{ijk}$，$\operatorname{pull}$ 读回聚合后的 $\bm{\theta}_i$。
增加服务器个数提高聚合带宽；分层同步先机架内再机架间。

\subsection{小结}

同步时间由拓扑决定。环同步适合 NVLink 稠密连接；
多机用参数服务器聚合 $\bm{g}_i=\sum_k\sum_j\bm{g}_{ijk}$，带宽不够就分层。

\subsection{练习}

核验：双向环能否减半延迟、计算尚未结束时异步发送对 $f$ 收敛的影响，以及单机失效时的容错。
